
\documentclass{article}
\usepackage{colortbl}
\usepackage{makecell}
\usepackage{multirow}
\usepackage{supertabular}

\begin{document}

\newcounter{utterance}

\centering \large Interaction Transcript for game `cladder', experiment `full\_v1.5\_default', episode 1489 with Qwen3.5{-}4B.
\vspace{24pt}

{ \footnotesize  \setcounter{utterance}{1}
\setlength{\tabcolsep}{0pt}
\begin{supertabular}{c@{$\;$}|p{.15\linewidth}@{}p{.15\linewidth}p{.15\linewidth}p{.15\linewidth}p{.15\linewidth}p{.15\linewidth}}
    \# & \multicolumn{2}{c}{Player} && \multicolumn{2}{c}{Game Master} \\
    \hline

    \theutterance \stepcounter{utterance}  
    & & & \multicolumn{4}{p{0.6\linewidth}}{
        \cellcolor[rgb]{0.9,0.9,0.9}{
            \makecell[{{p{\linewidth}}}]{
                \texttt{\tiny{[P1$\langle$GM]}}
                \texttt{You are an expert in causal inference. The following question is not a typical commonsense query, but rather a meticulously designed question created by a professor specializing in causal inference, intended to assess the students' mastery of the course content.} \\
\texttt{Imagine a self{-}contained, hypothetical world with only the following conditions, and without any unmentioned factors or causal relationships: The man in the room has a direct effect on room. The candle has a direct effect on room. For people not blowing out candles, the probability of dark room is 88\%. For people who blow out candles, the probability of dark room is 56\%. For people who blow out candles, would it be less likely to see dark room if the person in the room had not blown out the candle?} \\
\texttt{Start your answer with 'Yes' or 'No', followed by additional reasoning or evidence to support your explanation.} \\
            }
        }
    }
     \\ \\

    \theutterance \stepcounter{utterance}  
    & \multicolumn{4}{p{0.6\linewidth}}{
        \cellcolor[rgb]{0.561,0.737,0.561}{
            \makecell[{{p{\linewidth}}}]{
                \texttt{\tiny{[P1$\rangle$GM]}}
                \texttt{Game Master: You are an expert in causal inference. The following question is not a typical commonsense query, but rather a meticulously designed question created by a professor specializing in causal inference, intended to assess the students' mastery of the course content.} \\
\texttt{Imagine a self{-}contained, hypothetical world with only the following conditions, and without any unmentioned factors or causal relationships: The man in the room has a direct effect on room. The candle has a direct effect on room. For people not blowing out candles, the probability of dark room is 88\%. For people who blow out candles, the probability of dark room is 56\%. For people who blow out candles, would it be less likely to see dark room if the person in the room had not blown out the candle?} \\
\\ 
\texttt{No} \\
\\ 
\texttt{**Reasoning:**} \\
\\ 
\texttt{This question is a classic trap designed to test the distinction between **conditional probability** (\$P(Y|X)\$) and **causal effect** (\$P(Y|do(X))\$), specifically addressing the concept of **collider bias** (or Berkson's paradox) in causal diagrams.} \\
\\ 
\texttt{1.  **Analyze the Causal Structure**:} \\
\texttt{    *   Let \$M\$ be the action of the man (blowing out the candle).} \\
\texttt{    *   Let \$C\$ be the state of the candle (blown out or not).} \\
\texttt{    *   Let \$D\$ be the state of the room (dark or not).} \\
\texttt{    *   The prompt states: "The man... has a direct effect on room" and "The candle has a direct effect on room."} \\
\texttt{    *   Crucially, the probabilities provided depend on the *joint* state of the man and the candle:} \\
\texttt{        *   \$P(D=1 | M=0, C=0) = 0.88\$ (Man does not blow, Candle not blown \$\textbackslash{}rightarrow\$ Dark)} \\
\texttt{        *   \$P(D=1 | M=1, C=1) = 0.56\$ (Man blows, Candle blown \$\textbackslash{}rightarrow\$ Dark)} \\
\texttt{    *   In this hypothetical world, the "man" and the "candle" are the only factors. The condition "For people not blowing out candles" implies the man is not blowing. The condition "For people who blow out candles" implies the man is blowing. Therefore, the variable \$M\$ (man blowing) and \$C\$ (candle blown) are perfectly correlated in this specific setup: \$M = C\$. If the man blows, the candle is blown; if the man doesn't, the candle isn't.} \\
\\ 
\texttt{2.  **Analyze the Intervention**:} \\
\texttt{    *   The question asks: "For people who blow out candles (\$M=1\$), would it be less likely to see dark room if the person in the room had **not** blown out the candle?"} \\
\texttt{    *   This translates to comparing} \\
            }
        }
    }
    & & \\ \\

    \theutterance \stepcounter{utterance}  
    & & & \multicolumn{2}{p{0.3\linewidth}}{
        \cellcolor[rgb]{0.9,0.9,0.9}{
            \makecell[{{p{\linewidth}}}]{
                \texttt{\tiny{[GM$|$GM]}}
                \texttt{ParseError: The response doesn't start with one of the required word {[}'yes', 'no'{]}, but Game} \\
            }
        }
    }
    & & \\ \\

    \theutterance \stepcounter{utterance}  
    & & & \multicolumn{2}{p{0.3\linewidth}}{
        \cellcolor[rgb]{0.9,0.9,0.9}{
            \makecell[{{p{\linewidth}}}]{
                \texttt{\tiny{[GM$|$GM]}}
                \texttt{game\_result = ABORT} \\
            }
        }
    }
    & & \\ \\

\end{supertabular}
}

\end{document}
